\chapter{Boolean operations for deterministic transformers}

\section{Complement}\label{appendix:det-not}

Let $\TF$ be a transformer, we construct $\TF'$ that flips the decision token of $\TF$, i.e., if $\TF^*(\alpha) = \qaccept$, then $\TF'^*(\alpha) = \qreject$ and if $\TF^*(\alpha) = \qreject$, then $\TF'^*(\alpha) = \qaccept$.
This is the only change made to $\TF$, thus $\TF'^i(\sigma) = \TF^i(\sigma)$ for every $i$ when $\sigma$ is not a decision symbol.

This is done by swapping the rows associated with $\qaccept$ and $\qreject$ in the matrix of the $\OUTPUT$ layer of $\TF$.
Thus, when the largest probability is associated with $\qaccept$ in $\TF$, that same number will be associated with $\qreject$ in $\TF'$.

Formally, $\TF'$ is defined exactly the same as $\TF$ except in the $\OUTPUT$ layer, where the matrix is different.
Let $M$ be the matrix of the $\OUTPUT$ layer of $\TF$, let $i$ be the position of the token $\qaccept$ and $j$ the position of $\qreject$, then the matrix of the $\OUTPUT$ layer of $\TF'$ will be $M' = PM$ where $P$ is the matrix that flips the rows $i$ and $j$.

\[
    M'_{kl} = \begin{cases}
        1 & \text{ if } k = l \text{, } k \neq i \text{ and } k \neq j \\
        1 & \text{ if } k = j \text{ and } l = i \\
        1 & \text{ if } k = i \text{ and } l = j \\
        0 & \text{ otherwise}
    \end{cases}
\]

Since the $\OUTPUT$ matrix is the only difference between $\TF$ and $\TF'$, they have the same budget and embedding size as $\TF$, therefore, $\TF'$ belongs to the same class as $\TF$.

\medskip

We can define the family $\{\TF'_n\}_{n \in \mathbb{N}}$ where $\TF'_i$ is the negation of $\TF_i$ constructed as explained in this section.
Then $\{\TF'_n\}_{n \in \mathbb{N}}$ computes the complement of $\{\TF_n\}_{n \in \mathbb{N}}$.
Since for every $i \in \mathbb{N}$, $\TF'_i$ has the same budget and embedding size as $\TF_i$, $\bar{\gL}$ is in $\CoTtndn$.


\section{Disjunction}\label{appendix:det-or}

Let $\tfone{\gL}$ and  $\tftwo{\gL}$ be languages in $\CoTtndn$, then there is a family of transformers $\{(\tfone{\TF})_n\}_{n \in \mathbb{N}}$ (resp. $\{(\tftwo{\TF})_n\}_{n \in \mathbb{N}}$) of embedding size $\tfone{d}(n)$ (resp. $\tftwo{d}(n)$) $\in O(d(n))$ that after $\tfone{T}(n)$ (resp. $\tftwo{T}(n)$) $\in T(n)$ steps of $\CoT$ computes $\tfone{\gL}$ (resp. $\tftwo{\gL}$).

In this appendix, we construct a family $\{\TF_n\}_{n \in \mathbb{N}}$ that computes $\tfone{\gL} \cup \tftwo{\gL}$.
It does it in $\tfone{T}(n) + \tftwo{T}(n) + 1 \in O(T(n))$ steps of $\CoT$ and it has embedding size $O(d(n))$.

As mentioned in the  main text, $\TF$ first runs $\tfone{\TF}$, then $\tftwo{\TF}$, and lastly, computes the disjunction of the outputs.

\medskip
First, observe that for every input $\alpha$, the tape of a transformer, at the end of the computation, ends up looking as $\alpha \beta q$, where $q$ is the decision token printed by the transformer and $\beta$ ($\in \Sigma^*$) is the sequence of intermediate tokens printed during the chain of thought.
We denote by $\alpha \beta \tfone{q}$ the content of the tape of $\tfone{\TF}$ at the end of the computation when the input word is $\alpha$.
Similarly, we denote by $\alpha \gamma \tftwo{q}$ the content of the tape of $\tftwo{\TF}$ at the end of the computation when the input word is $\alpha$.
The tape of $\TF$, computing the disjunction of $\tfone{\TF}$ and $\tftwo{\TF}$, at the end of the computation will have content $\alpha \beta \tfone{q} \gamma \tftwo{q} q$, where $q$ is the decision token representing the disjunction of $\tfone{q}$ and $\tftwo{q}$.


\medskip

Let us construct $\TF$ from $\tfone{\TF}$ and $\tftwo{\TF}$.
Since we cannot change the body of the transformer during the iterations, we are required to use some tricks.

First, we grow the dimension of the embedding for running the body of both transformers at the same time over the same vectors but independently.
In the final layers we choose which probability distribution will go to the output layer, this makes $\TF$ print the token generated by $\tfone{\TF}$ or $\tftwo{\TF}$.
As always, we assume that the transformers are normalized.

We use the first half of the embedding for the transformer $\tfone{\TF}$ and the second for the transformer $\tftwo{\TF}$.
We have some extra coordinates for flags at the end, but we will ignore them for now.

\[
\TE(\sigma)_i = \begin{cases}
    \tfone{\TE}(\sigma)_i                       &\text{ if } i \leq \tfone{d}(|\alpha|) \\
    \tftwo{\TE}(\sigma)_{i-\tfone{d}(|\alpha|)} &\text{ if } \tfone{d}(|\alpha|) < i \leq \tftwo{d}(|\alpha|) \\
    0                                           &\text{ if } \tftwo{d}(|\alpha|) < i\\
\end{cases}
\]\[
\PE(n)_i = \begin{cases}
    \tfone{\PE}(n)_i                            &\text{ if } i \leq \tfone{d}(|\alpha|) \\
    \text{shifted}\tftwo{\PE}(n)_{i-\tfone{d}(|\alpha|)}\footnotemark      &\text{ if } \tfone{d}(|\alpha|) < i \leq \tftwo{d}(|\alpha|) \\
    \text{flags}                                &\text{ if } \tftwo{d}(|\alpha|) < i\\
\end{cases}
\]\footnotetext{Defined later in the construction}


The transformer $\TF$ has one layer for each layer of $\tfone{\TF}$ and each of $\tftwo{\TF}$.
The first layers correspond to the body of $\tfone{\TF}$.
They are constructed ignoring the half of the embedding storing the coordinates used by $\tftwo{\TF}$.
The second half of layers are the other way around, they compute the layers of $\tftwo{\TF}$ and ignore the $\tfone{\TF}$ coordinates.
Then, all the matrices in the layers of the first part of the body of $\TF$ look like $V_1$ and the ones in the second layers will look like $V_2$ where $\tfone{W}$ and $\tftwo{W}$ are the matrices of the corresponding layers in $\tfone{\TF}$ and $\tftwo{\TF}$.

\begin{align*}
    &V_1 = \left(\begin{matrix}
        \tfone{W} & 0 \\
        0  & 0 \\
    \end{matrix}\right)
    &V_2 = \left(\begin{matrix}
        0 & 0 \\
        0  & \tftwo{W} \\
    \end{matrix}\right)
\end{align*}

It is important to notice that the values of the coordinates corresponding to one transformer do not affect the values of the other.
Thus, if the tape of the transformer has content $\gamma$, then at the end of the body of $\TF$, the last vector would look like $(x, y, \text{flags})$ where $x$ is the last vector at the end of the body of $\tfone{\TF}$ when its tape content is $\gamma$, and $y$ the one at the end of $\tftwo{\TF}$ when its tape content is $\gamma$.
Therefore, if we always choose to transform said vector into $(x, 0, \dots, 0)$, then $\TF$ will produce the same tape as $\tfone{\TF}$.
This works for generating the first $\tfone{T}(|\alpha|)$ tokens which are $\beta \tfone{q}$.
Thus, we would like to apply this transformation only to the vectors generating them.
Therefore, we need to flag them with the position encoding.
Since $h_{|\alpha| + i}$ carries the distribution to output in the $(i+1)$-th step, i.e., the $(|\alpha| + i + 1)$-th token of the tape, the vectors that we need to flag will be $h_{|\alpha|}, h_{|\alpha|+1}, \dots, h_{|\alpha|+\tfone{T}(|\alpha|)-1}$.

\smallskip

After running $\tfone{T}(|\alpha|)$ steps of this construction, the tape will be $\alpha \beta \tfone{q}$.
At this point, we want to start generating $\gamma$, i.e., the tape of $\tftwo{\TF}$.
Just changing what we choose before the $\OUTPUT$ layer does not work, since what the construction will produce would be $\tftwo{\TF}(\alpha \beta \tfone{q})$, and we need $\tftwo{\TF}(\alpha)$.

\medskip

Now we show how to run $\tftwo{\TF}$ in this construction ignoring the tokens produced by $\tfone{\TF}$.
We make $\tftwo{\TF}$ not attend to the tokens in positions $|\alpha| + 1$ to $\tfone{T}(|\alpha|)$ (this can be done using the primitive presented in the Section \ref{ignore-vecs}).
We also need for them to be invisible to the portion of the position encoding corresponding to $\tftwo{\TF}$, thus we "shift" it in the following way.

\[
\PE(n)_i = \begin{cases}
    \tfone{\PE}(n)_i                                        &\text{ if } i \leq \tfone{d}(|\alpha|) \\
    \text{shifted}\tftwo{\PE}(n)_{i-\tfone{d}(|\alpha|)}    &\text{ if } \tfone{d}(|\alpha|) < i \leq \tftwo{d}(|\alpha|) \\
    \text{flags}                                            &\text{ if } \tftwo{d}(|\alpha|) < i\\
\end{cases}
\]

\[
\text{shifted}\tftwo{\PE}(n) = \begin{cases}
        \tftwo{\PE}(n)                          &\text{ if } n \leq |\alpha|             \\
        0                                       &\text{ if } |\alpha| < n < \tfone{T}(|\alpha|) \\
        \tftwo{\PE}(n - \tfone{T}(|\alpha|))    &\text{ if } \tfone{T}(|\alpha|) < n                   \\
    \end{cases}
\]

The first case is for the positions occupied by the input, we do not want to change their encoding.
The second case is for the positions occupied by the tokens generated by $\tfone{\TF}$, we do not care about them.
The third case is for the tokens that $\tftwo{\TF}$ will generate.
It should believe that they are in their natural positions, i.e., shifted $\tfone{T}(|\alpha|)$ positions left (the positions occupied by the tokens generated by $\tfone{\TF}$).
We also need to change the flag used for choosing the distribution in the vectors at positions after $|\alpha| + \tfone{T}(|\alpha|)$ so $\TF$ chooses the distribution corresponding to $\tftwo{\TF}$ instead of $\tfone{\TF}$.

At this point, the construction almost produces the tape of $\tfone{\TF}$ followed by the tape of $\tftwo{\TF}$ except for one token.
We need to do something special for the first token generated by $\tftwo{\TF}$.
The distribution corresponding to it is in the vector $h_{|\alpha|}$ but the last vector of the first run computing $\tftwo{\TF}$ tokens is $h_{|\alpha| + \tfone{T}(|\alpha|)}$, the one coming from the last token of $\tfone{\TF}$.
Therefore, before choosing the distribution, we need to copy the vector $h_{|\alpha|}$ to $h_{|\alpha| + \tfone{T}(|\alpha|)}$.
Notice that nothing more is needed, since the attention mechanism only attends to vectors to the right, so $h_{|\alpha|}$ is not affected by the vectors coming from the tokens produced by $\tfone{\TF}$.
This is done with the operation presented in Section \ref{copy-vector}.


With this last fix, $\TF$ generates the tape of $\tfone{\TF}$ followed by the tape of $\tftwo{\TF}$.
We now show how to compute $q$, i.e., the disjunction of $\tfone{q}$ and $\tftwo{q}$.
This is easily done making $h_{|\alpha| + \tfone{T}(|\alpha|) + \tftwo{T}(\alpha)}$ only attend to itself and to $h_{|\alpha| + \tfone{T}(|\alpha|)}$, i.e., to the vectors coming from $\tfone{q}$ and $\tftwo{q}$.


We add a flag in the token embedding of $\TF$ to recognize the decision tokens of $\tfone{\TF}$ and $\tftwo{\TF}$.
The flag is in the coordinate $f_c$.

\[
    \TE(\sigma)_{f_c} = \begin{cases}
        1 & \text{ if } \sigma = \qaccept \\
        0 & \text{ otherwise } \\
    \end{cases}
\]

Then, adding $h_{|\alpha| + \tfone{T}(|\alpha|)}$ to $h_{|\alpha| + \tfone{T}(|\alpha|) + \tftwo{T}(\alpha)}$, we have $1$ or $2$ in the coordinate $f_c$ iff $\tfone{\TF}$ or $\tftwo{\TF}$ accepted $\alpha$.
Now we make all of $h_{|\alpha| + \tfone{T}(|\alpha|) + \tftwo{T}(\alpha)}$ $0$, except the $f_c$ coordinate (using Section \ref{const-fun}).

Before going to the $\OUTPUT$ layer we move the value of $f_c$ to the coordinate that will project to the probability of $\qaccept$ and add some number larger than $0$ but smaller than $1$ to the coordinate that will define the probability of $\qreject$.
Therefore, the token with the maximum probability will be $\qaccept$ iff $\tfone{\TF}$ or $\tftwo{\TF}$ accepted $\alpha$, and it will be $\qreject$ in the other case.

Then, $\TF$ computes the union of the languages of $\tfone{\TF}$ and $\tftwo{\TF}$.
The number of steps of $\CoT$ that it needs is exactly the ones made by $\tfone{\TF}$ plus the ones made by $\tftwo{\TF}$ plus $1$ extra for taking the disjunction, thus it has the same asymptotic cost as the maximum of $\tfone{\TF}$ and $\tftwo{\TF}$.
The embedding dimension of $\TF$ is the same as the dimension of $\tfone{\TF}$ plus the dimension of $\tftwo{\TF}$ plus some constant number of flags, thus in $O(\cdot)$ notation it is the same as the maximum of $\tfone{\TF}$ and $\tftwo{\TF}$.

Therefore, the language computed by the family of transformers $\{\TF_n\}_{n \in \mathbb{N}}$ is in $\CoTtndn$.
Thus, $\gL_1 \cup \gL_2 \in \CoTtndn$.

